
Our work connects three research areas: parameter-efficient fine-tuning, sub-quadratic architectures, and transfer learning. We position our contribution at their intersection, revealing architectural constraints that prior work has not systematically explored.

\subsubsection{Parameter-Efficient Fine-Tuning}

Low-rank adaptation (LoRA) introduced a method for fine-tuning large models by injecting low-rank weight updates into linear layers, reducing trainable parameters while matching full fine-tuning performance. Subsequent work explored alternative strategies: prefix tuning prepends learnable tokens to inputs, prompt tuning optimizes continuous prompts, and adapters insert bottleneck layers. These methods share a common assumption: the base model architecture already possesses capabilities needed for the target task.

All established PEFT research focuses on transformer architectures---BERT, GPT-2, T5, and variants. Transformers' bidirectional self-attention enables flexible information flow suitable for diverse tasks. This architectural versatility allowed PEFT literature to treat the base model as a black box, focusing on which parameters to adapt rather than whether adaptation is possible. Our work challenges this assumption by testing PEFT applicability to architecturally distinct model families.

\subsubsection{State-Space Models and Sub-Quadratic Architectures}

State-space models emerged as efficient alternatives to transformers, replacing quadratic attention with linear-time recurrent dynamics. Structured state-space models (S4) introduced parameterization enabling stable training of deep SSM layers, achieving competitive performance on long-range benchmarks. Mamba extended S4 with selective state-space mechanisms and hardware-aware implementations, demonstrating strong results on language modeling and generation tasks.

SSM research has primarily evaluated generative capabilities---language modeling perplexity, long-context understanding, and sequence generation. Classification tasks, especially those requiring bidirectional reasoning, remain underexplored. Mamba's design inherits causal constraints from language modeling: the model processes sequences left-to-right, with each token's representation depending only on preceding context. This architectural choice optimizes generation but may impose fundamental limitations on tasks requiring symmetric comparison.

Other sub-quadratic architectures---linear attention, RWKV, and Hyena---share similar efficiency motivations but differ in implementation. Our focus on Mamba as a representative causal state-space model provides a testbed for understanding how architectural constraints affect PEFT applicability.

\subsubsection{Transfer Learning and Architectural Compatibility}

Classical transfer learning research established that pretrained representations accelerate downstream task learning. Recent work explored domain shift, task similarity, and few-shot adaptation. These studies assume architectural compatibility: source and target tasks must be solvable by the same model architecture.

Zero-shot evaluation has been used primarily to assess pretrained model quality rather than as an architectural compatibility test. Our work reframes zero-shot performance as a critical gating mechanism: below-random accuracy signals architectural mismatch that PEFT cannot overcome. This shifts zero-shot evaluation from a benchmark metric to a prerequisite validation step.

\subsubsection{Positioning Our Contribution}

Prior PEFT research demonstrates how to efficiently adapt transformers but does not address when adaptation is possible across architectures. SSM research validates efficiency but has not systematically tested architectural limitations on classification tasks. Transfer learning assumes compatibility but lacks principled methods to validate it a priori.

We contribute systematic evaluation of PEFT applicability across architecture families, revealing that causal state-space models fail bidirectional tasks at zero-shot evaluation---a failure pattern that fine-tuning cannot overcome. Our task-architecture compatibility framework establishes architectural validation as a prerequisite for PEFT, preventing wasted compute on structurally impossible experiments. As the field explores diverse efficient architectures, understanding these constraints becomes essential for practitioners deploying PEFT methods beyond transformers.
